% Nouveau EGA V (cote 25): archive cote 25, site page 75 (PDF page 76).
% typescript. Draft compared: none.
% AI-transcribed by Claude Opus 5.5 (claude-opus-5-5, adaptive thinking, effort high), Message Batches API, batch ega25b (msgbatch_01A7n5YuXubUPNK1zsepi8RG), 2026-11-13; not yet checked by eye.
\scanpage{25/75}
\mspage{70}

20.13. \uline{Morphismes élémentaires et théorème de M. Artin}.

\uline{Définition} 13.1. Un morphisme $f:X\to Y$ de préschémas est dit ``morphisme
 élémentaire'' si X est  Y-isomorphe à un préschéma de la forme
$X' - Z$ , où $X'\to X$ est un morphisme lisse, projectif, à fibres géomé-%
triques connexes de dimension 1, et où Z est un sous-préschéma fermé
de X',  tel que le morphisme $Z\to Y$ soit étale, surjectif, et
de degré constant. Un morphisme est dit  poly-élémen-%
taire s'il est un composé de morphismes élémentaires. Un préschéma X
sur un corps k est dit polyélementaire (sur k), si le morphisme structu-%
ral $X\to \mathrm{Spec}(k)$ est poly-élémentaire.

\uline{Théorème} 13.2. (M.Artin) Soient X un préschéma géom. irréd. sur un corps k parfait infini, x un point lisse
de X, alors x admet un système fondamental de voisinages ouverts
polyélémentaires.

Quitte à remplacer X par un voisinage donné de x,  il revient au
 même de montrer qu'il existe un voisinage ouvert élémentaire de x dans X.
  Raisonnant par récurrence sur
la dimension n de X, on  est ramené à montrer que si $n>0$ , il existe
un voisinage ouvert U de x et un morphisme élémentaire $f: U\to V$,
V étant un schéma lisse sur k (nécessairement géom. irréd et de dim n-1).
 (Le cas n=0 est évidemment trivial, étant donné qu' alors
X est k-isomorphe à Spec(k), qui est polyélémentaire sur Spec(k), étant
entendu dans 13.1. qu'on n'exc\hand{l}ut pas le composé de la famille vide de
morphismes - on aurait dû le signaler sous une forme ou une autre dans
13.1. ). La nécessité de supposer k parfait au préalable apparait déjà
dans le cas n=1: alors on prend pour X' le modèle projectif normal
canonique du corps des fonctions K de X (cf Chap II) par.7), le fait que

\mspage{75}
